% Experiments. Contract: ../references/section-playbook.md
% 600-1000 words in four subsections: setup, main results, distribution, ablations.
\section{Experiments}
\label{sec:experiments}

\subsection{Experimental setup}
\label{sec:setup}

\textbf{Data and evaluation.}
All runs use MNIST with the upstream four-way partition into training,
validation, disagreement, and test sets. The \evalsem\ are frozen as described in
\cref{sec:setting}: $\delta = 0.05$, a 1000-point bound-inversion grid, and
10{,}000 Monte Carlo samples for the disagreement estimate. These values are
identical for every run below and were fixed before the comparison matrix was
registered.

\textbf{Scale.}
Runs use a reduced probe scale of 12 epochs and a maximum compression size of
120, against the upstream configuration's 200 epochs and 20{,}000. Scale is a
free parameter of this study, but it is applied identically on both sides of
every comparison, so a change in scale cannot produce a difference between
methods. \Cref{tab:scale} reports what changes when the compression budget is
doubled.

\textbf{Compared configurations.}
\tightcolorbox{ACCompare}{\textbf{\methodcontrol.}} The unmodified pipeline,
re-run and recorded in this environment.
\tightcolorbox{ACPrimary}{\textbf{Selection methods.}} \methodmax\ and
\methodrandom, differing only in which examples the \surrogate\ is built from,
as described in \cref{sec:methods}.

\textbf{Seeds and repetition.}
5 \splitseed{}s $\{1,2,3,4,42\}$, 3 trajectory seeds per split and method, giving
45 preserved primary runs. Statistics follow \cref{eq:persplit} and
\cref{eq:contrast}; the reporting threshold of \cref{eq:decision} was fixed at
$\kappa = 1.0$ before execution.

\textbf{Environment and failure handling.}
All runs executed on a single GPU under one dependency lock, recorded in each run
manifest and reproduced in \cref{sec:appendix-config}. No run in the reported
matrix crashed. Crashed runs, had there been any, would be preserved with their
manifests and counted rather than silently retried or dropped from the
denominator.

\subsection{Main results}
\label{sec:main-results}

% Mandatory main results table. Contract: references/table-grammar.md, pattern 1.
% Every cell is filled from runs/aggregate__<method>.json.
\begin{table}[t!]
  \caption{\textbf{\certbound\ (\down) under the frozen protocol.}
    \swatch{ACCompare} is the \methodcontrol,
    \swatch{ACPrimary} are the compared selection methods.
    Mean is over 5 \splitseed{}s, each aggregating 3 trajectory seeds under
    identical \evalsem. Std.\ is the standard deviation of the per-split means
    and Range is their max\,$-$\,min; both describe variation \emph{between}
    splits, not within one. The paired contrast is computed against the
    \methodcontrol\ within split and trajectory index, following
    \cref{eq:contrast}; ``--'' marks the row a contrast is taken against.}
  \label{tab:main-results}
  \vspace{-0.2cm}
  \centering
  {\fontsize{8pt}{11pt}\selectfont
   \resizebox{\linewidth}{!}{%
     \begin{NiceTabular}{l c c c c c c c}[colortbl-like]
       \toprule[1.2pt]
         & \begin{tabular}{c}\textbf{Split}\\\textbf{seeds}\end{tabular}
         & \begin{tabular}{c}\textbf{Traj.}\\\textbf{per split}\end{tabular}
         & \textbf{Mean}
         & \textbf{Std.}
         & \textbf{Range}
         & \begin{tabular}{c}\textbf{Max traj.}\\\textbf{std.}\end{tabular}
         & \begin{tabular}{c}\textbf{Paired}\\\textbf{contrast}\end{tabular} \\
       \midrule[1.2pt]\midrule
       \rowcolor{ACCompare}\methodcontrol
         & 5 & 3 & 0.409 & 0.007 & 0.019 & 0.019 & -- \\
       \midrule
       \rowcolor{ACPrimary}\methodmax
         & 5 & 3 & 0.398 & 0.009 & 0.024 & 0.023 & $-0.011$ \\
       \rowcolor{ACPrimary}\methodrandom
         & 5 & 3 & 0.387 & 0.013 & 0.030 & 0.032 & $-0.022$ \\
       \bottomrule
     \end{NiceTabular}
   }
  }
\end{table}


\textbf{Does the selection method change the \certbound?}
\Cref{tab:main-results} gives the aggregate view. Both selection methods improve
on the \methodcontrol: \methodmax\ by $0.011$ and \methodrandom\ by $0.022$ in
paired within-split contrast. Between the two selection methods the paired
contrast is $0.011 \pm 0.005$ in favour of \methodrandom. The largest within-split
trajectory spread anywhere in the matrix is $0.032$, so by \cref{eq:decision} the
between-method contrast does not clear the noise it is measured against. The
supported statement is that the two methods are indistinguishable at this scale,
with a direction that favours \methodrandom\ but is not established.

\textbf{Is the ordering stable across splits?}
% Mandatory secondary table. Contract: references/table-grammar.md, pattern 2.
\begin{table}[t!]
  \caption{\textbf{Per-split \certbound\ and paired contrast (\down).} Each
    column is one \splitseed. Every entry is the mean over the same three
    trajectory seeds, so the difference row is an exact pairing. A positive
    difference favors \methodrandom. The sign reverses at \splitseed~42, which
    is why \cref{tab:main-results} reports a contrast and not a ranking.}
  \label{tab:paired}
  \vspace{-0.2cm}
  \centering
  {\fontsize{8pt}{11pt}\selectfont
     \begin{NiceTabular}{l c c c c c}[colortbl-like]
       \toprule[1pt]
       \textbf{\splitseed} & \textbf{1} & \textbf{2} & \textbf{3} & \textbf{4} & \textbf{42} \\
       \midrule
       \methodcontrol & 0.412 & 0.405 & 0.418 & 0.399 & 0.410 \\
       \methodmax     & 0.401 & 0.395 & 0.412 & 0.388 & 0.394 \\
       \rowcolor{ACPrimary}\methodrandom & 0.379 & 0.384 & 0.398 & 0.371 & 0.401 \\
       \midrule
       Difference     & 0.022 & 0.011 & 0.014 & 0.017 & $-0.007$ \\
       \bottomrule
     \end{NiceTabular}
  }
\end{table}

\Cref{tab:paired} reports the same comparison split by split. The contrast
favours \methodrandom\ at four of the five splits and reverses at \splitseed~42,
where \methodmax\ is lower by $0.007$. Under the design rule of
\cref{sec:factors}, a ranking requires consistency across splits, so no ranking
is reported. A mean contrast computed over splits whose individual contrasts
disagree in sign describes the sample, not the method.

Notice also that the between-split range grows monotonically with how aggressively
the method selects: $0.019$ for the \methodcontrol, $0.024$ for \methodmax, and
$0.030$ for \methodrandom. Whatever a selection criterion does to the mean, it
also changes how much the certificate moves when the split changes, and the
latter effect is the larger one here.

\subsection{Trajectory variation}
\label{sec:trajectory}

\IfFileExists{data/bounds.dat}{%
  % Tier B results figure: pgfplots reading a data file written by the analysis
% script. Guarded with \IfFileExists so that, before the data exists, the float
% does not appear at all (Tier D of references/figure-ladder.md).
%
% source-data: runs/aggregate__{control,maxerr,random}.json -> paper/data/bounds.dat
% generator:   scripts/make_paper_data.py --figure distribution
\begin{figure}[!t]
  \centering
  \begin{tikzpicture}
    \begin{axis}[acplot,
      xlabel={\splitseed}, ylabel={\certbound\ $\downarrow$},
      xtick=data, symbolic x coords={1,2,3,4,42},
      enlarge x limits=0.12,
    ]
      \addplot[name path=lo,draw=none,forget plot] table[x=seed,y=rlo] {data/bounds.dat};
      \addplot[name path=hi,draw=none,forget plot] table[x=seed,y=rhi] {data/bounds.dat};
      \addplot[ACBlue!14,forget plot] fill between[of=lo and hi];
      \addplot+[ACMuted, mark=triangle*, dashed]
        table[x=seed,y=control] {data/bounds.dat};
      \addplot+[ACCompare!55!black, mark=square*]
        table[x=seed,y=maxerr] {data/bounds.dat};
      \addplot+[ACBlue, mark=*]
        table[x=seed,y=random] {data/bounds.dat};
      \legend{\methodcontrol, \methodmax, \methodrandom}
    \end{axis}
  \end{tikzpicture}
  \caption{\textbf{Per-split \certbound\ and trajectory spread.} Each marker is
    the mean over 3 trajectory seeds at a fixed \splitseed, under identical
    \evalsem; the shaded band is the trajectory range of \methodrandom. The band
    contains the \methodmax\ mean at every split, so the ordering visible in the
    means is not separable from trajectory noise. The comparison is read across
    splits, not at a single split: the ordering reverses at \splitseed~42.}
  \label{fig:distribution}
\end{figure}
%
  \Cref{fig:distribution} shows the per-split means for all three
  configurations, with the trajectory range of \methodrandom\ as a shaded band.
  The band is wide enough at every split to contain the corresponding
  \methodmax\ mean, which is the same conclusion \cref{eq:decision} reaches
  arithmetically.%
}{}

Within-split trajectory standard deviations range from $0.011$ to $0.032$ across
the matrix, and are systematically larger for \methodrandom\ than for the
\methodcontrol. Against that, the between-method contrast of $0.011$ is small. A
comparison at this scale that used a single trajectory per configuration would
have had roughly a one-in-three chance of reporting the opposite ordering, which
is the practical reason the matrix is specified with repetition rather than
breadth. Where an individual trajectory is discussed, it is discussed as an
example: a single optimization path is an observation about that run, not an
ablation result and not evidence about a policy.

\subsection{Ablations}
\label{sec:ablations}

% Paired ablation tables. Contract: references/table-grammar.md, pattern 3.
\begin{table}[t!]
  \centering
  \begin{minipage}{0.48\textwidth}
    \centering
    \captionsetup{justification=centering,singlelinecheck=false}
    \caption{\textbf{Effect of trajectory count (\down).}}
    \vspace{-0.1cm}
    \label{tab:trajectory-count}
    {\fontsize{8pt}{11pt}\selectfont
     \begin{NiceTabular}{l ccc}[colortbl-like]
       \toprule[1.2pt]
       \textbf{Statistic} & \textbf{1} & \textbf{2} & \textbf{3} \\
       \midrule[1.2pt]
       Mean \certbound       & 0.362 & 0.376 & 0.387 \\
       Between-split std.    & 0.013 & 0.013 & 0.013 \\
       Mean paired contrast  & 0.014 & 0.012 & 0.011 \\
       \bottomrule[1.2pt]
     \end{NiceTabular}%
    }
  \end{minipage}
  \hfill
  \begin{minipage}{0.44\textwidth}
    \centering
    \captionsetup{justification=centering,singlelinecheck=false}
    \caption{\textbf{Effect of experiment scale (\down).}}
    \vspace{-0.1cm}
    \label{tab:scale}
    {\fontsize{8pt}{11pt}\selectfont
     \begin{NiceTabular}{l ccc}[colortbl-like]
       \toprule[1.2pt]
       \textbf{Setting} & \textbf{Epochs} & \textbf{Comp.} & \textbf{Bound} \\
       \midrule[1.2pt]
       Reduced      &  12 &   120 & 0.387 \\
       Intermediate &  12 &   240 & 0.381 \\
       Original     & 200 & 20000 & --    \\
       \bottomrule[1.2pt]
     \end{NiceTabular}%
    }
  \end{minipage}
\end{table}


\Cref{tab:trajectory-count} varies the number of trajectory seeds aggregated per
split. The estimated mean drifts upward as trajectories are added, from $0.362$
at one trajectory to $0.387$ at three, while the between-split standard deviation
is unchanged at $0.013$ and the mean paired contrast shrinks from $0.014$ to
$0.011$. A one-trajectory estimate is therefore both optimistic about the bound
and generous about the contrast, in the direction that would most easily produce
a positive-sounding result. Three trajectories is where the contrast stops
moving, not where it clears the threshold of \cref{eq:decision}.

\Cref{tab:scale} compares compression budgets. Doubling the budget from 120 to
240 lowers the bound for \methodrandom\ by $0.006$, an effect of the same order
as the between-method contrast. The comparison at the upstream scale of 200
epochs and 20{,}000 was not run, and no value is reported for it.
